%------------------------------------------------------------------------------%
\begin{question}
  Mostre que
  \[
    f(x, y) =
    \begin{cases}
      \dfrac{2xy}{x^2+y^2}, & (x, y) \neq (0, 0) \\
      \quad\; 0,            & (x, y) =    (0, 0)
    \end{cases}
  \]
  é contínua em todos os pontos, exceto na origem
\end{question}

%------------------------------------------------------------------------------%
\begin{resolution}{Solução}
  Consideramos primeiro o caso \((x, y) \neq (0, 0)\)

  \pause
  \vspace{\baselineskip}
  \(2xy\) é uma função contínua no plano

  \pause
  \(x^2+y^2\) é uma função contínua no plano

  \pause
  então \(\frac{2xy}{x^2+y^2}\)
  \pause
  é contínua para todo \((x, y) \neq (0, 0)\)
\end{resolution}

%------------------------------------------------------------------------------%
\begin{resolution}{Solução}
  Em \((x, y) = (0, 0)\) temos que aplicar a definição de continuidade

  \pause
  \vspace{0.5\baselineskip}
  \begin{enumerate}
    \setlength\itemsep{1.5em}
    \item \(f\) é definida em \((0, 0)\)
    \item \(\lim_{(x, y)\to(0, 0)} f(x, y)\) existe
    \item \(\lim_{(x, y)\to(0, 0)} f(x, y) = f(0, 0)\)
  \end{enumerate}

  \pause
  \vspace{\baselineskip}
  A condição 1 vale, pois \(f(0, 0) = 0\)
\end{resolution}

%------------------------------------------------------------------------------%
\begin{resolution}{Solução}
  Verificando a condição 2, percebemos que o limite não existe

  \pause
  Se fizermos \(y=mx\),
  \pause
  para \(m\neq 0\) e \(x\neq 0\), temos
  \pause
  \[
    f(x, y)\evalat{y=mx}{}
    \pause
    = \frac{2xy}{x^2+y^2}\evalat{y=mx}{}
    \pause
    = \frac{2mx^2}{x^2+m^2x^2}
    \pause
    = \frac{2mx^2}{\left(1+m^2\right)x^2}
    \pause
    = \frac{2m}{1+m^2}
  \]

  \pause
  Para cada valor de $m$, a função $f$ se aproxima da origem com um valor diferente.

  \pause
  Portanto, o limite não existe e a função não é contínua na origem
\end{resolution}

%------------------------------------------------------------------------------%
